\documentclass[12pt,a4paper]{report}

%% ── Packages ──────────────────────────────────────────────────────────────────
\usepackage[a4paper, margin=2.5cm]{geometry}
\usepackage{fontenc}
\usepackage{inputenc}
\usepackage{microtype}
\usepackage{lmodern}
\usepackage{xcolor}
\usepackage{graphicx}
\usepackage{float}
\usepackage{booktabs}
\usepackage{longtable}
\usepackage{tabularx}
\usepackage{multirow}
\usepackage{array}
\usepackage{listings}
\usepackage{caption}
\usepackage{subcaption}
\usepackage{hyperref}
\usepackage{fancyhdr}
\usepackage{titlesec}
\usepackage{tcolorbox}
\usepackage{enumitem}
\usepackage{amsmath}
\tcbuselibrary{skins,breakable}

%% ── Colors ────────────────────────────────────────────────────────────────────
\definecolor{primary}{HTML}{1E3A5F}
\definecolor{accent}{HTML}{2563EB}
\definecolor{lightblue}{HTML}{EFF6FF}
\definecolor{success}{HTML}{059669}
\definecolor{warning}{HTML}{D97706}
\definecolor{codeblue}{HTML}{0369A1}
\definecolor{codegray}{HTML}{F1F5F9}
\definecolor{codered}{HTML}{BE123C}

%% ── Hyperref Setup ────────────────────────────────────────────────────────────
\hypersetup{
  colorlinks=true, linkcolor=accent, urlcolor=accent, citecolor=accent,
  pdftitle={FinSight -- Intelligent Personal Finance Analytics Platform (Part 2)},
  pdfauthor={Aayank Singh}
}

%% ── Header / Footer ───────────────────────────────────────────────────────────
\pagestyle{fancy}
\fancyhf{}
\fancyhead[L]{\small\textcolor{primary}{\textbf{FinSight} -- Intelligent Personal Finance Analytics Platform}}
\fancyhead[R]{\small\textcolor{primary}{SPE Major Project}}
\fancyfoot[C]{\thepage}
\renewcommand{\headrulewidth}{0.4pt}

%% ── Section Formatting ────────────────────────────────────────────────────────
\titleformat{\chapter}[display]
  {\normalfont\huge\bfseries\color{primary}}
  {\chaptertitlename\ \thechapter}{20pt}{\Huge}
\titleformat{\section}{\normalfont\Large\bfseries\color{primary}}{\thesection}{1em}{}
\titleformat{\subsection}{\normalfont\large\bfseries\color{accent}}{\thesubsection}{1em}{}
\titleformat{\subsubsection}{\normalfont\normalsize\bfseries\color{primary}}{\thesubsubsection}{1em}{}

%% ── Code Listings ─────────────────────────────────────────────────────────────
\lstset{
  basicstyle=\ttfamily\footnotesize,
  backgroundcolor=\color{codegray},
  frame=single, framerule=0pt, rulecolor=\color{codegray},
  breaklines=true, breakatwhitespace=true,
  keywordstyle=\color{codeblue}\bfseries,
  commentstyle=\color{success}\itshape,
  stringstyle=\color{codered},
  numberstyle=\tiny\color{gray},
  numbers=left, stepnumber=1,
  tabsize=2, showspaces=false, showstringspaces=false,
  captionpos=b
}

%% ── Custom Boxes ──────────────────────────────────────────────────────────────
\newtcolorbox{infobox}[1][]{
  enhanced, colback=lightblue, colframe=accent,
  fonttitle=\bfseries\color{white}, title=#1,
  boxrule=0.5pt, arc=4pt, left=6pt, right=6pt, top=6pt, bottom=6pt
}

\begin{document}

% Continue chapter numbering from Part 1
\setcounter{chapter}{4}

%% ════════════════════════════════════════════════════════════════════════════
\chapter{Microservices Specification \& DB Schemas}
%% ════════════════════════════════════════════════════════════════════════════

\section{Service Breakdown}

The FinSight platform is decomposed into 14 distinct core backend microservices. The separation between Java Spring Boot and Python FastAPI is strictly enforced based on the nature of workload (user-facing transactional vs. intensive machine learning inference).

\subsection{Java Spring Boot Microservices}

\subsubsection{1. API Gateway Service (\texttt{api-gateway-service})}
Acts as the single front door for all client traffic. Built on Spring Cloud Gateway, it implements global filters for authorization. It intercepts incoming HTTP requests, extracts the RS256 JWT, validates its signature against the public key, and enforces rate limiting (100 requests/min/user) using a Redis token bucket algorithm.

\subsubsection{2. User Service (\texttt{user-service})}
Manages user onboarding, authentication, profile metadata, and session lifecycle. Supports email/password registration with bcrypt hashing (cost factor 12) and Google OAuth 2.0.
\begin{itemize}[leftmargin=1.5em, noitemsep]
  \item \textbf{DB Schema Table (\texttt{user\_credentials}):} \texttt{id} (UUID), \texttt{email} (VARCHAR 320, UNIQUE), \texttt{password\_hash} (VARCHAR 255), \texttt{full\_name}, \texttt{city}, \texttt{age\_group}, \texttt{monthly\_income}, \texttt{role} (VARCHAR 20), \texttt{deactivated\_at}.
  \item \textbf{Token Tracking Tables:} \texttt{refresh\_token\_sessions} (\texttt{id}, \texttt{token\_jti}, \texttt{user\_email}, \texttt{expires\_at}) and \texttt{revoked\_access\_tokens}.
\end{itemize}

\subsubsection{3. Ingestion Service (\texttt{ingestion-service})}
Orchestrates the ingestion of bank statement files (CSV, XLS, PDF). It stores uploaded file metadata and asynchronous parse job documents in MongoDB. Once parsed into canonical rows, it publishes \texttt{transactions.ingested} events to Kafka.

\subsubsection{4. Transaction Service (\texttt{transaction-service})}
The core ledger of the platform. Consumes ingestion events, performs deduplication using a cryptographic hash (\texttt{content\_hash} generated from date + amount + description), and manages categories.
\begin{itemize}[leftmargin=1.5em, noitemsep]
  \item \textbf{DB Schema Table (\texttt{transactions}):} \texttt{id} (UUID), \texttt{content\_hash} (VARCHAR 64, UNIQUE), \texttt{owner\_email}, \texttt{account\_id}, \texttt{job\_id}, \texttt{source\_file\_name}, \texttt{source\_bank}, \texttt{occurred\_at}, \texttt{raw\_description}, \texttt{normalized\_merchant}, \texttt{category\_id}, \texttt{amount} (NUMERIC 15,2), \texttt{type} (DEBIT/CREDIT), \texttt{currency}, \texttt{tags} (TEXT[]), \texttt{is\_anomaly}.
\end{itemize}

\subsubsection{5. Budget Service (\texttt{budget-service})}
Enables users to set monthly spending caps per category. Consumes transaction events from Kafka to update running totals in real time. If spend breaches 80\% or 100\% of a budget limit, it publishes events to \texttt{budget.alerts}.
\begin{itemize}[leftmargin=1.5em, noitemsep]
  \item \textbf{DB Schema Table (\texttt{budgets}):} \texttt{id} (UUID), \texttt{owner\_email}, \texttt{category\_id}, \texttt{category\_name}, \texttt{limit\_amount}, \texttt{current\_spend}, \texttt{month\_year} (YYYY-MM).
\end{itemize}

\subsubsection{6. Notification Service (\texttt{notification-service})}
Processes \texttt{budget.alerts} and \texttt{anomaly.alerts} from Kafka. It persists alerts in a notification inbox and pushes real-time updates over WebSocket/STOMP to active client sessions.
\begin{itemize}[leftmargin=1.5em, noitemsep]
  \item \textbf{DB Schema Table (\texttt{notifications}):} \texttt{id} (UUID), \texttt{owner\_email}, \texttt{title}, \texttt{type}, \texttt{message}, \texttt{is\_read}, \texttt{created\_at}.
  \item \textbf{DB Schema Table (\texttt{notification\_preferences}):} \texttt{owner\_email}, \texttt{type}, \texttt{enabled}.
\end{itemize}

\subsubsection{7. Chat Service (\texttt{chat-service})}
Provides a conversational AI interface using the Gemini API. It intercepts natural language questions, queries internal microservices (transaction, budget, forecasting) via REST WebClient for grounding context, and synthesizes personalized responses.
\begin{itemize}[leftmargin=1.5em, noitemsep]
  \item \textbf{DB Schema Tables:} \texttt{chat\_sessions} (\texttt{id}, \texttt{owner\_email}), \texttt{chat\_messages} (\texttt{id}, \texttt{session\_id}, \texttt{role}, \texttt{content}, \texttt{intent}), \texttt{chat\_ratings} (\texttt{message\_id}, \texttt{rating}).
\end{itemize}

\subsubsection{8. Admin Service (\texttt{admin-service})}
Provides administrative oversight. Exposes endpoints for user deactivation, platform telemetry aggregation, and dynamic ML threshold adjustment.
\begin{itemize}[leftmargin=1.5em, noitemsep]
  \item \textbf{DB Schema Table (\texttt{admin\_config}):} \texttt{config\_key} (PRIMARY KEY), \texttt{config\_value}, \texttt{description}. Stores parameters like \texttt{anomaly.sensitivity} and \texttt{stress.weight.*}.
\end{itemize}

\subsection{Python FastAPI ML Microservices}

\subsubsection{9. Parsing ML Service (\texttt{parsing-ml-service})}
Utilizes NLP and heuristic column mapping to intelligently extract transaction records from unstructured or multi-column bank statements without requiring hardcoded templates.

\subsubsection{10. Categorization Service (\texttt{categorization-service})}
Employs an XGBoost classifier trained on merchant names and transaction descriptions to classify transactions into standard financial categories. Listens to Kafka for unclassified transactions and publishes categorized outcomes.

\subsubsection{11. Anomaly Detection Service (\texttt{anomaly-detection-service})}
Runs a PyTorch-based LSTM autoencoder. For each user, it builds an individual spending DNA baseline. It scores incoming transactions in near real-time (target latency $<$ 200ms) and flags anomalous behavior.

\subsubsection{12. Forecasting Service (\texttt{forecasting-service})}
Uses time-series forecasting algorithms (LSTM/ARIMA) over historical transaction data to predict next month's spending across categories, enabling proactive budget vs. prediction comparisons.

\subsubsection{13. Stress Score Service (\texttt{stress-score-service})}
Computes a monthly composite financial stress score (0--100) combining spend-to-income ratios, savings rates, recurring payment burden, and budget adherence.

\subsubsection{14. Simulation Service (\texttt{simulation-service})}
Powers the interactive "What-If" simulator in the frontend. Evaluates hypothetical spending modifications in real-time by chaining forecasting and stress scoring models without mutating production state.

\subsubsection{15. Model Training Service (\texttt{model-training-service})}
Manages scheduled and admin-triggered model retraining. Computes KL divergence to detect data drift and integrates with MLflow for model tracking and artifact versioning.

%% ════════════════════════════════════════════════════════════════════════════
\chapter{Cloud-Native Infrastructure \& DevOps Automation}
%% ════════════════════════════════════════════════════════════════════════════

\section{Docker Containerization}

Every microservice in FinSight is fully containerized. A multi-stage build approach is adopted for Java Spring Boot services using Maven to compile and package the application, resulting in lean runtime images based on Eclipse Temurin 21 Alpine. Python services use \texttt{uvicorn} with lean Python 3.11 slim base images.

\begin{figure}[H]
  \centering
  \begin{tcolorbox}[enhanced, colback=gray!10, colframe=gray!50, width=0.85\linewidth,
      height=5cm, arc=6pt, valign=center, halign=center]
    \textcolor{gray!70}{\large\textit{[SCREENSHOT PLACEHOLDER: Docker Compose Infrastructure Setup]}}
  \end{tcolorbox}
  \caption{Docker Compose Multi-Container Deployment}
\end{figure}

\section{HashiCorp Vault Integration}

To achieve a zero-trust security posture, no plaintext passwords or API keys exist in source code or configuration files. HashiCorp Vault is deployed as the central secrets engine.

\begin{itemize}[leftmargin=1.5em]
  \item \textbf{Secret Engine Path:} Secrets are stored under \texttt{secret/data/finsight}.
  \item \textbf{Kubernetes Injection:} In Kubernetes, a bridge mechanism (External Secrets Operator / init-containers) populates Kubernetes Secrets (\texttt{finsight-secrets}) directly from Vault.
  \item \textbf{Runtime Binding:} Deployments inject these secrets as environment variables into pods (e.g., \texttt{DB\_PASSWORD}, \texttt{JWT\_PRIVATE\_KEY}, \texttt{GEMINI\_API\_KEY}).
\end{itemize}

\section{Kubernetes Architecture \& Configuration}

FinSight is deployed into a dedicated Kubernetes namespace (\texttt{finsight}). The architecture leverages Deployments, StatefulSets, Services, ConfigMaps, and Secrets.

\begin{figure}[H]
  \centering
  \begin{tcolorbox}[enhanced, colback=gray!10, colframe=gray!50, width=0.9\linewidth,
      height=6cm, arc=6pt, valign=center, halign=center]
    \textcolor{gray!70}{\large\textit{[SCREENSHOT PLACEHOLDER: Kubernetes Cluster Workloads \& Pods]}}
  \end{tcolorbox}
  \caption{Kubernetes Pods, Services, and StatefulSets in \texttt{finsight} Namespace}
\end{figure}

\subsection{Stateful Data Storage (PostgreSQL)}
To ensure reliable persistence and predictable DNS resolution, PostgreSQL is deployed as a \texttt{StatefulSet} with a Headless Service (\texttt{clusterIP: None}). It attaches a \texttt{PersistentVolumeClaim} requesting 5Gi of ReadWriteOnce storage.

\subsection{Probes \& Rolling Updates}
All Deployments configure strict liveness and readiness probes pointing to \texttt{/health} or Actuator endpoints. To satisfy the zero-downtime requirement, the Deployment strategy is explicitly configured:
\begin{lstlisting}[language=YAML, caption=Kubernetes RollingUpdate Strategy]
strategy:
  type: RollingUpdate
  rollingUpdate:
    maxUnavailable: 0
    maxSurge: 1
\end{lstlisting}

\section{Horizontal Pod Autoscaling (HPA)}

Machine learning inference (especially LSTM scoring) is highly CPU-intensive. To maintain SLA under traffic spikes without over-provisioning infrastructure, \texttt{HorizontalPodAutoscaler} resources are bound to \texttt{anomaly-detection-service} and \texttt{categorization-service}.

\begin{lstlisting}[language=YAML, caption=HPA Configuration for ML Inference Services]
apiVersion: autoscaling/v2
kind: HorizontalPodAutoscaler
metadata:
  name: anomaly-detection-service-hpa
  namespace: finsight
spec:
  scaleTargetRef:
    apiVersion: apps/v1
    kind: Deployment
    name: anomaly-detection-service
  minReplicas: 2
  maxReplicas: 8
  metrics:
    - type: Resource
      resource:
        name: cpu
        target:
          type: Utilization
          averageUtilization: 70
\end{lstlisting}
When average CPU utilization across pods exceeds 70\%, Kubernetes automatically scales the deployment up to 8 replicas.

\section{Ansible Automation}

Infrastructure provisioning and application deployment across environments are automated using Ansible. The setup mandates four distinct roles:
\begin{enumerate}[leftmargin=1.5em]
  \item \textbf{\texttt{common}:} Prepares base OS, installs prerequisite libraries, and configures Docker.
  \item \textbf{\texttt{kafka}:} Deploys and configures Apache Kafka and Zookeeper clusters.
  \item \textbf{\texttt{kubernetes}:} Applies Kubernetes namespace, ConfigMaps, Secrets, and manifest files idempotently.
  \item \textbf{\texttt{monitoring}:} Deploys Prometheus, Grafana, and ELK Stack agents.
\end{enumerate}

\section{Jenkins CI/CD Pipeline}

Continuous integration and deployment are driven by a robust \texttt{Jenkinsfile} triggered automatically via GitHub webhooks on push to the \texttt{main} branch.

\begin{figure}[H]
  \centering
  \begin{tcolorbox}[enhanced, colback=gray!10, colframe=gray!50, width=0.9\linewidth,
      height=5cm, arc=6pt, valign=center, halign=center]
    \textcolor{gray!70}{\large\textit{[SCREENSHOT PLACEHOLDER: Jenkins CI/CD Pipeline Execution]}}
  \end{tcolorbox}
  \caption{Jenkins 7-Stage Automated CI/CD Pipeline}
\end{figure}

The pipeline executes the following sequential stages:
\begin{enumerate}[leftmargin=1.5em, noitemsep]
  \item \textbf{Checkout:} Pulls the latest source code from GitHub.
  \item \textbf{Build Backend:} Patches internal service discovery URLs and compiles Java microservices using Maven (\texttt{mvn package}).
  \item \textbf{Build Frontend:} Compiles the React/TypeScript frontend via Vite (\texttt{npm run build}).
  \item \textbf{Test:} Executes unit and integration test suites, generating JUnit/Surefire XML reports.
  \item \textbf{Build \& Push Docker Images:} Builds Docker images for all 12 services and pushes them to Docker Hub tagged as \texttt{latest}.
  \item \textbf{Deploy with Ansible:} Executes the Ansible deploy playbook (\texttt{deploy.yml}) to pull images and update container orchestrators.
  \item \textbf{Post-Build Cleanup:} Prunes dangling images to free worker node disk space.
\end{enumerate}

%% ════════════════════════════════════════════════════════════════════════════
\chapter{Observability \& Logging Pipeline}
%% ════════════════════════════════════════════════════════════════════════════

\section{ELK Stack Architecture}

A comprehensive centralized logging pipeline is established using Elasticsearch, Logstash, and Kibana. 

\begin{figure}[H]
  \centering
  \begin{tcolorbox}[enhanced, colback=gray!10, colframe=gray!50, width=0.9\linewidth,
      height=6cm, arc=6pt, valign=center, halign=center]
    \textcolor{gray!70}{\large\textit{[SCREENSHOT PLACEHOLDER: Kibana Log Analytics Dashboard]}}
  \end{tcolorbox}
  \caption{Kibana Centralized Telemetry and Log Analytics Dashboard}
\end{figure}

\begin{itemize}[leftmargin=1.5em]
  \item \textbf{Log Emission:} All Spring Boot and FastAPI microservices emit structured JSON logs to TCP port 5044.
  \item \textbf{Logstash Aggregation:} Logstash ingests JSON streams, applies grok/filtering rules, and indexes records into Elasticsearch.
  \item \textbf{Kibana Visualizations:} Pre-built Kibana dashboards monitor real-time transaction ingestion rates, ML inference latencies, anomaly trigger volumes, and per-service error rates.
  \item \textbf{Distributed Tracing:} A \texttt{correlation\_id} is injected into request headers at the API Gateway and propagated across HTTP/Kafka hops for end-to-end trace reconstruction.
\end{itemize}

\section{Prometheus \& Grafana Telemetry}

System and JVM telemetry are scraped by Prometheus and visualized in Grafana.
\begin{itemize}[leftmargin=1.5em]
  \item \textbf{Metrics Exposition:} Spring Boot microservices utilize Micrometer to expose Prometheus metrics at \texttt{/actuator/prometheus}.
  \item \textbf{Grafana Dashboards:} Real-time monitoring panels visualize CPU throttling, memory consumption, JVM heap usage, and HTTP request throughput per pod.
\end{itemize}

%% ════════════════════════════════════════════════════════════════════════════
\chapter{Innovative Solutions Implemented}
%% ════════════════════════════════════════════════════════════════════════════

\section{Personalized Anomaly Detection (LSTM Autoencoder)}
Standard financial platforms employ static rule-based thresholds (e.g., flagging transactions over ₹50,000). FinSight utilizes a PyTorch LSTM autoencoder that learns the temporal sequence of an individual user's spending habits (merchant frequency, typical transaction times, amount distribution). A ₹500 transaction at 11 PM at a petrol pump is considered normal for a sales executive but flagged as highly anomalous for a student.

\section{Composite Financial Stress Score}
FinSight computes an advanced monthly financial health metric (0--100) aggregating five critical financial vectors:
\begin{equation}
\text{Stress Score} = w_1\left(\frac{\text{Spend}}{\text{Income}}\right) + w_2(\text{Savings Rate}) + w_3(\text{EMI Burden}) + w_4(\text{Discretionary Growth}) + w_5(\text{Budget Breach})
\end{equation}
Where weights $w_i$ are dynamically configurable via the Admin Service. Changes greater than 10 points month-on-month trigger proactive alerts.

\section{What-If Financial Simulator}
An interactive simulation engine allowing users to manipulate sliders (e.g., "Reduce dining spend by ₹2,000" or "Add new EMI of ₹5,000"). The frontend calls the \texttt{simulation-service}, which dynamically evaluates the impact on next month's cash flow and stress score without modifying live ledger state.

\begin{figure}[H]
  \centering
  \begin{tcolorbox}[enhanced, colback=gray!10, colframe=gray!50, width=0.85\linewidth,
      height=5cm, arc=6pt, valign=center, halign=center]
    \textcolor{gray!70}{\large\textit{[SCREENSHOT PLACEHOLDER: What-If Simulator UI]}}
  \end{tcolorbox}
  \caption{What-If Financial Simulator Interface}
\end{figure}

\section{Spending Mood Correlation}
Users optionally log daily mood check-ins (Happy, Stressed, Sad). Over time, machine learning correlates emotional states with transaction categories, generating behavioral insights such as: \textit{"You spend 60\% more on online food delivery on days rated as Stressed."}

\section{Contextual AI Chat Assistant}
Rather than providing static FAQ responses, the AI Chat Assistant utilizes Gemini LLM grounded in real user data. When a user queries \textit{"Why did I overspend in March?"}, the service orchestrates internal calls to the anomaly service and budget service, synthesizing a precise, data-backed natural language explanation.

%% ════════════════════════════════════════════════════════════════════════════
\chapter{Conclusion \& Roadmap}
%% ════════════════════════════════════════════════════════════════════════════

FinSight successfully demonstrates the convergence of modern personal finance analytics, machine learning, and enterprise DevOps automation. By decoupling services via Kafka, securing secrets with Vault, ensuring observability through ELK/Prometheus, and implementing zero-downtime Kubernetes deployments, the platform establishes a highly scalable and robust architecture. Future development roadmap includes multi-bank account aggregation via Account Aggregator (AA) frameworks and automated tax harvesting optimization.

\end{document}
